\chapter{Kidney Stones}
\index{Kidney Stones}

\section*{Definition and Overview}
Kidney stones are hard deposits of minerals and salts that form inside the kidneys when the urine becomes concentrated and the substances that form stones crystallise. They range from tiny grains to large stones, and they may cause no symptoms while sitting in the kidney. When a stone moves into the ureter, it can cause sudden, severe pain in the flank radiating to the groin, blood in the urine, and, if it blocks urine flow, infection and kidney damage. Most stones pass on their own, but larger stones need treatment such as lithotripsy or surgery. Drinking plenty of fluid and dietary changes are the main ways to prevent new stones.

\section*{Epidemiology}
Around one in ten people will have a kidney stone at some time, and rates are rising, partly because of increasing rates of obesity, diabetes, and dietary changes. Men are affected about twice as often as women, and stones are most common between ages 30 and 60. Recurrence is common: around half of people who have had one stone will have another within five to ten years. Calcium stones are the most common type, followed by uric acid, struvite, and cystine stones.

\section*{Aetiology and Pathophysiology}
Stones form when the urine is supersaturated with stone-forming substances, so that crystals form and grow instead of being flushed out. The most common is calcium oxalate, promoted by low fluid intake, high salt and animal protein, and conditions that increase calcium in the urine. Uric acid stones form when the urine is acidic, as in gout and with high animal protein intake. Struvite stones follow urinary infections. Cystine stones occur in an inherited condition. Dehydration, hot climates, obesity, diabetes, and a family history all increase risk. A stone forms in the kidney and usually passes down the ureter; most that are under 5 millimetres pass spontaneously, while larger stones become stuck.

\section*{Clinical Features}
\begin{itemize}
    \item No symptoms while the stone stays in the kidney
    \item Sudden, severe colicky pain in the flank or back, radiating to the groin and genitals
    \item Waves of pain lasting 20 to 60 minutes
    \item Blood in the urine, which may be visible or found on testing
    \item Nausea and vomiting
    \item Pain or burning with urination, urgency, and frequency
    \item Fever and chills, which indicate an infected or obstructed stone
    \item Difficulty passing urine if the stone blocks both ureters or the outlet
    \item Restlessness and inability to find a comfortable position
\end{itemize}

\section*{Differential Diagnosis}
Other causes of severe flank and abdominal pain must be excluded.
\begin{itemize}
    \item \textbf{Kidney infection:} fever with flank pain
    \item \textbf{Urinary tract infection:} burning and frequency
    \item \textbf{Bowel conditions:} appendicitis, diverticulitis, and intestinal obstruction
    \item \textbf{Gallbladder disease:} right upper abdominal pain
    \item \textbf{Musculoskeletal pain:} back strain and disc disease
    \item \textbf{Gynaecological conditions:} ovarian cysts and ectopic pregnancy
    \item \textbf{Pancreatitis:} upper abdominal pain radiating to the back
\end{itemize}

\section*{When to Seek Medical Care}
Anyone with sudden severe flank pain, blood in the urine, or fever should be assessed. Urgent care is needed for pain with fever and shaking (which may indicate an infected, obstructed stone), inability to pass urine, or severe vomiting. People who have passed a stone should have it analysed, and those with recurrent stones should be investigated for the cause.

\textbf{Contact your local emergency services for:}
\begin{itemize}
    \item Severe flank pain with fever, shaking, or confusion, which may indicate sepsis
    \item Inability to pass urine
    \item Vomiting with an inability to keep fluids down
    \item Blood in the urine with severe pain
    \item Collapse, dizziness, or severe illness
\end{itemize}

\section*{Diagnosis and Investigations}
\begin{itemize}
    \item \textbf{History and examination:} the pattern of pain and risk factors
    \item \textbf{Urine tests:} blood, infection, crystals, and pH
    \item \textbf{Blood tests:} kidney function, calcium, uric acid, and electrolytes
    \item \textbf{CT scan:} the best imaging for stones, showing size and position
    \item \textbf{Ultrasound:} an alternative that avoids radiation
    \item \textbf{Stone analysis:} of any stone passed, to guide prevention
    \item \textbf{Metabolic testing:} 24-hour urine collection for recurrent stones
\end{itemize}

\section*{Management}
\begin{itemize}
    \item \textbf{Fluids and pain relief:} the mainstay for small stones, with most passing within days to weeks
    \item \textbf{Medicines:} medicines that help stones pass and relax the ureter
    \item \textbf{Lithotripsy:} shock waves to break stones into passable fragments
    \item \textbf{Ureteroscopy:} a scope passed into the ureter to remove or break the stone
    \item \textbf{Keyhole or open surgery:} for very large stones
    \item \textbf{Stenting:} a tube to bypass an obstructing stone
    \item \textbf{Urgent treatment:} for infected or completely obstructing stones
    \item \textbf{Prevention:} fluids, dietary change, and medicines for recurrent stones
\end{itemize}

\section*{Complications}
\begin{itemize}
    \item Obstruction with kidney damage if a stone blocks urine flow
    \item Infection and sepsis in an obstructed kidney
    \item Recurrent stones
    \item Chronic pain and repeated presentations
    \item Complications of treatment
    \item Reduced kidney function with large or repeated stones
\end{itemize}

\section*{Prognosis and Outcomes}
Most kidney stones pass without treatment, and those that need intervention are treated successfully with lithotripsy or surgery in the great majority of cases. The main challenge is recurrence, which affects half of people within a decade. Increasing fluid intake alone reduces recurrence by about half, and dietary change and medicines help those with specific causes. With prevention, most people with kidney stones avoid further episodes.

\section*{Prevention and Self-Care}
\begin{itemize}
    \item Drink 2 to 3 litres of fluid a day, enough that the urine is pale
    \item Limit salt and reduce animal protein
    \item Eat calcium-rich foods, as dietary calcium reduces stone risk
    \item Avoid excessive vitamin C and oxalate supplements if prone
    \item Limit foods high in oxalate if advised
    \item Have stones analysed and seek metabolic testing after recurrence
    \item Maintain a healthy weight
    \item Take prescribed preventive medicines
    \item Seek early care for fever or severe pain
\end{itemize}

\section*{Sources and Further Reading}
The following reputable sources provide further information for healthcare students and patients seeking reliable, current advice about kidney stones.
\begin{itemize}
    \item Kidney Health Australia. \textit{Kidney stones.} https://kidney.org.au/
    \item Urological Society of Australia and New Zealand. \textit{Kidney stone information.} https://www.usanz.org.au/
    \item Healthdirect Australia. \textit{Kidney stones.} https://www.healthdirect.gov.au/
    \item Better Health Channel. \textit{Kidney stones.} https://www.betterhealth.vic.gov.au/
    \item NHS. \textit{Kidney stones.} https://www.nhs.uk/conditions/kidney-stones/
    \item National Kidney Foundation. \textit{Kidney stones.} https://www.kidney.org/
\end{itemize}
